\documentclass{article}
\usepackage{graphicx}

\usepackage[utf8]{inputenc} % Para aceitar acentos
\usepackage[T1]{fontenc}    % Para renderizar fontes com acentos corretamente
\usepackage[brazil]{babel}  % Para português (hifenização, datas, etc.)
\usepackage[a4paper, margin=2.5cm]{geometry} % Ajusta a largura da folha
\usepackage[ruled,vlined, linesnumbered]{algorithm2e}% Escrever pseudocódigo
\usepackage{amsmath}
\usepackage{amssymb}
\usepackage{hyperref}
\usepackage{xcolor}
\newcommand{\blue}[1]{\textcolor{blue}{#1}}
\newcommand{\red}[1]{\textcolor{red}{#1}}

\setlength{\parindent}{0pt}
\setlength{\parskip}{4pt}

\title{O problema da compra mínima: formulações e heurísticas.}
\author{Antônio Erick Freitas Ferreira - 322980\\
        Rafael C. S. Schouery \\
        \small Instituto de Computação - Unicamp - São Paulo - Brasil
}

\begin{document}
\maketitle

    \begin{abstract}

        \textbf{AND:} Relações comerciais estão presentes no nosso cotidiano. Um vendedor expõe seus produtos à venda e um comprador que esteja disposto a desembolsar o valor pode adquiri-lo. Com um mercado cada vez mais competitivo, o desafio do vendedor passa a ser como precificar seus produtos.

        \textbf{BUT:} No entanto, surge um dilema econômico: preços mais elevados aumentam a receita unitária, mas podem reduzir o conjunto de compradores dispostos a adquirir determinado produto; por outro lado, preços mais baixos ampliam o número de vendas possíveis, porém diminuem o ganho obtido em cada transação. Assim, do ponto de vista do vendedor, o problema consiste em definir preços para um conjunto de itens de forma a equilibrar essas forças antagônicas e maximizar a receita total obtida.

        Com o objetivo de maximizar a receita do vendedor, define-se o \textit{problema da compra mínima} como segue. Seja $I$ um conjunto de itens, $B$ um conjunto de compradores e $V$ uma matriz em que cada elemento $v_{ib}$ representa o valor que o comprador $b \in B$ está disposto a pagar pelo item $i \in I$. O problema consiste em determinar um conjunto de preços $P$, em que cada elemento $p_i$ corresponde ao preço atribuído ao item $i$. Além disso, o comportamento do comprador é conhecido: ele comprará um único produto mais barato viável, caso contrário, ele não comprará nada.

        \textbf{THEREFORE:} Portanto, este trabalho tem como objetivo propor melhorias às formulações matemáticas existentes, com o intuito de fortalecer os modelos de Programação Linear Inteira, bem como desenvolver heurísticas eficientes para a resolução de instâncias de grande porte do problema.

        \vspace{0.5cm}
        
        \center \textbf{Palavras-chave:} Otimização combinatória, precificação, heurística.
    \end{abstract}

    Este trabalho tem como objetivo investigar possíveis melhorias nas modelagens e implementações já existentes para o problema, bem como propor heurísticas baseadas em algoritmos genéticos para resolução de problemas de grande porte em tempo hábil.
\end{document}